%% LyX 1.5.5 created this file.  For more info, see http://www.lyx.org/.
%% Do not edit unless you really know what you are doing.
\documentclass[english]{article}
\usepackage[T1]{fontenc}
\usepackage[latin9]{inputenc}
\usepackage{esint}
\usepackage{babel}

\begin{document}

\title{Three temperatures distribution function model}


\author{by Yves Peysson (CEA/IRFM)}

\maketitle
The three temperatures model has been introduced for simulating an
anistropic tail in the electron momentum distribution function driven
by the quasi-electrostatic wave at the Lower Hybrid frequency in tokamak.
It is a very crude model which was used to characterize the fast electron
bremsstrahlung during this non-inductive regimes, which extends up
to high photon energies corresponding to hard X rays. 

With the emergence of refined model based on the use of Fokker-Planck
solvers, the three temperatures model is almost obsolete, and its
application is restricted to benchmark non-thermal bremsstrahlung
codes designed for analysing corresponding diagnostics.

Let define the 3-D distribution function $f\left(\psi,p,\xi\right)$,
where $\psi$ is the poloidal magnetic flux coordinate (radial position),
$p$ the momentum value and $\xi=p_{\parallel}/p$ the pitch-angle
relative to the local magnetic field direction. Here, $f$ is considered
to be homogeneous on a magnetic flux surface, and the trapped electron
population is neglected. According to the three temperatures model,\begin{equation}
f\left(\psi,p,\xi\right)=\left(1-\lambda\right)f_{M}\left(\psi,p,\xi\right)+\lambda f_{3T}\left(\psi,p,\xi\right)\label{eq:distribution}\end{equation}
where $f_{M}\left(\psi,p,\xi\right)$ is the thermal bulk, and $f_{3T}\left(\psi,p,\xi\right)$
is the non-thermal part. By definition, both $f_{M}$ and $f_{3T}$
are normalized to unity, so that $f$ automatically satisfies this
condition. The relative fraction of fast electrons is given by the
parameter $\lambda$ which is usually very small of the ordre of $1\times10^{-3}.$


\section{The Maxwellian distribution function}

The thermal distribution function is expressed in its general relativistic
form\begin{equation}
f_{M}\left(\psi,p,\xi\right)\propto\alpha_{M}\exp\left[-\frac{p^{2}}{\left(1+\gamma\right)\Theta}\right]=\alpha_{M}\exp\left[-\frac{\gamma-1}{\Theta}\right]\label{eq:Maxwellian}\end{equation}
where $\Theta\left(\psi\right)=T_{e}\left(\psi\right)/m_{e}c^{2}$
is the ratio of the local electron temperature $T_{e}\left(\psi\right)$
to the electron rest mass energy $m_{e}c^{2}$ while the relativistic
Lorentz correction factor is defined as $\gamma^{2}=p^{2}+1.$ Here,
$\xi$ is the cosine of the pitch angle. By symmetry, the Maxwellian
distribution function is independant of $\xi$. By definition, $f_{M}$
is normalized to unity in the interval $p\in\left[0,p_{min}\right[$,
where $p_{min}$ is a given upper cut-off limit for the Maxwellian.
The parameter $\alpha_{M}$$ $ is then given by the integral \begin{equation}
4\pi\alpha_{M}\int_{0}^{p_{min}}p^{2}\exp\left[-\frac{\gamma-1}{\Theta}\right]dp=1\end{equation}
Recalling that $\gamma d\gamma=pdp,$ \begin{equation}
\alpha_{M}=\frac{1}{4\pi\exp\left(\Theta^{-1}\right)\Xi_{p}\left(\Theta^{-1},p_{min}\right)}\end{equation}
where \begin{equation}
\Xi_{p}\left(\Theta^{-1},p_{min}\right)=\int_{1}^{\sqrt{1+p_{min}^{2}}}\gamma\sqrt{\gamma^{2}-1}\exp\left[-\gamma\Theta^{-1}\right]d\gamma\end{equation}


By integrating by parts\begin{eqnarray}
\Xi_{p}\left(\Theta^{-1},p_{min}\right) & = & \frac{p_{min}^{3}}{3}\exp\left[-\sqrt{1+p_{min}^{2}}\Theta^{-1}\right]\nonumber \\
 &  & +\frac{\Theta^{-1}}{3}\int_{1}^{\sqrt{1+p_{min}^{2}}}\left(\gamma^{2}-1\right)^{3/2}\exp\left[-\gamma\Theta^{-1}\right]d\gamma\label{eq:}\end{eqnarray}
and in the limit $p_{min}\rightarrow+\infty$, \begin{eqnarray}
\Xi_{p}\left(\Theta^{-1}\right) & = & \frac{\Theta^{-1}}{3}\int_{1}^{\infty}\left(\gamma^{2}-1\right)^{3/2}\exp\left[-\gamma\Theta^{-1}\right]d\gamma\\
 & = & \frac{4}{3}\frac{\left(3/2\right)!}{\sqrt{\pi}}\frac{K_{2}\left(\Theta^{-1}\right)}{\Theta^{-1}}\end{eqnarray}
where $K_{2}\left(z\right)$ is the modified Bessel function of order
$2$. 

For $\Theta\ll1,$ using the large argument asymptotic development
of $K_{2}\left(\Theta^{-1}\right)\approx\sqrt{\pi/2}\sqrt{\Theta}\exp(-\Theta^{-1})+O\left(\Theta^{-3/2}\right),$
the usual expression \begin{eqnarray}
f_{M}\left(\psi,p,\xi\right) & \simeq & \frac{1}{\left[2\pi\Theta\right]^{3/2}}\exp\left[-\frac{\gamma-1}{\Theta}\right]\\
 & = & \frac{1}{\left[2\pi\Theta\right]^{3/2}}\exp\left[-\frac{p^{2}}{\left(1+\gamma\right)\Theta}\right]\end{eqnarray}
is well recovered and \begin{equation}
\alpha_{M}\simeq\left[2\pi\Theta\right]^{-3/2}\label{eq:}\end{equation}


In the LUKE code, the momentum $p$ here expressed in relativistic
units $m_{e}c$ is normalized to the thermal reference value $p_{th}^{\dagger}=m_{e}v_{th}^{\dagger},$\begin{equation}
\overline{p}=p/p_{th}^{\dag}\end{equation}
 and consequently \begin{equation}
p_{th}^{\dagger}/m_{e}c\approx v_{th}^{\dagger}/c=\beta_{th}^{\dagger}\end{equation}
 since thermal electrons are only weakly relativistic. The well known
relativistic Lorentz correction factor $\gamma$ is then simply given
by the relation \begin{equation}
\gamma=\sqrt{p^{2}+1}=\sqrt{\overline{p}^{2}\beta_{th}^{\dagger2}+1}\end{equation}
 and in the non-relativistic limit, i.e. the condition $\overline{p}^{2}\beta_{th}^{\dag2}\ll1$
or $\gamma\approx1$ holds.

Since in relativistic units, the total energy is linked to the relativistic
momentum by the expression, \begin{equation}
\left(E_{c}+1\right)^{2}=\gamma\end{equation}
 it is straightforward to express the kinetic energy $E_{c}$ as a
function of $\overline{p}$ in units of electron rest mass energy
$m_{e}c^{2}$\begin{equation}
E_{c}=m_{e}c^{2}\left(\sqrt{\overline{p}^{2}\beta_{th}^{\dagger2}+1}-1\right)\end{equation}


Finally, concerning the normalization of the electron velocity $v$,
one has \begin{eqnarray}
-1\ll v/c & = & p/\left(\gamma m_{e}c\right)\\
 & = & \overline{p}p_{th}^{\dagger}/\left(\gamma m_{e}c\right)\\
 & = & \overline{p}\beta_{th}^{\dagger}/\gamma\end{eqnarray}
 and using $\overline{v}=v/v_{th}^{\dagger},$ it comes \begin{equation}
\overline{v}v_{th}^{\dagger}/c=\overline{p}\beta_{th}^{\dagger}/\gamma\end{equation}
 or \begin{equation}
\overline{v}=\overline{p}/\gamma\end{equation}
 with \begin{equation}
v_{th}^{\dagger}/c=\beta_{th}^{\dagger}\end{equation}


Consequently, the approximate expression of $f_{M}$ is only valid
when $\gamma\simeq1$, or in an equivalent form \begin{equation}
\overline{p}^{2}\beta_{th}^{\dagger2}\ll1\label{weakrelativist}\end{equation}
Since $\overline{p}$ may be as large as $30$ in numerical calculations,
in order to correctly describe momentum dynamics of the fastest electrons,
it results that \begin{equation}
\beta_{th}^{\dagger}=\Theta^{\dagger2}\ll1/30\end{equation}
or \begin{equation}
T_{e}^{\dagger}\ll511\sqrt{1/30}\approx100keV\end{equation}
since $\beta_{th}^{\dagger}=\sqrt{T_{e}^{\dagger}/m_{e}c^{2}}$ and
$\Theta^{\dagger}=T_{e}^{\dagger}/m_{e}c^{2}$. Usually, this condition
is always well satisfied even in the core of tokamak plasmas, where
$T_{e}^{\dagger}$ never exceeds $20keV.$ In normalized units,

\begin{eqnarray}
f_{M}\left(\psi,\overline{p},\xi\right) & = & \frac{1}{\left[2\pi\Theta\left(\psi\right)\right]^{3/2}}\exp\left[-\frac{\overline{p}^{2}\beta_{th}^{\dagger2}}{\left(1+\gamma\right)\overline{T}_{e}\left(\psi\right)\beta_{th}^{\dagger2}}\right]\\
 & = & \frac{1}{\left[2\pi\overline{T}_{e}\left(\psi\right)\beta_{th}^{\dagger2}\right]^{3/2}}\exp\left[-\frac{\overline{p}^{2}}{\left(1+\gamma\right)\overline{T}_{e}\left(\psi\right)}\right]\end{eqnarray}
using the relation $\Theta=\overline{\Theta}\Theta^{\dagger}$, with
$ $$\overline{\Theta}=\overline{T}_{e}$. Then, it turns out that
\begin{equation}
f_{M}\left(\psi,\overline{p},\xi\right)=\frac{1}{\beta_{th}^{\dagger3}}\overline{f}_{M}\left(\psi,\overline{p},\xi\right)=\frac{1}{p_{th}^{\dagger3}}\overline{f}_{M}\left(\psi,\overline{p},\xi\right)\end{equation}
 since $p_{th}^{\dagger}=\beta_{th}^{\dagger}$, with \begin{equation}
\overline{f}_{M}\left(\psi,\overline{p},\xi\right)\approx\frac{1}{\left[2\pi\overline{T}_{e}\left(\psi\right)\right]^{3/2}}\exp\left[-\frac{\overline{p}^{2}}{\left(1+\gamma\right)\overline{T}_{e}\left(\psi\right)}\right]\end{equation}
and the modified normalisation coefficient becomes \begin{equation}
\overline{\alpha}_{M}\simeq\left[2\pi\overline{T}_{e}\left(\psi\right)\right]^{-3/2}\label{eq:}\end{equation}


The Maxwellian distribution function be expressed in an alternative
form, useful for calculating interpolation between full and half-grids,
\begin{equation}
\overline{f}_{M}\left(\psi,\overline{p},\xi\right)\approx\frac{1}{\left[2\pi\overline{T}_{e}\left(\psi\right)\right]^{3/2}}\exp\left[-\frac{\gamma-1}{\overline{T}_{e}\left(\psi\right)\beta_{th}^{\dagger2}}\right]\end{equation}


One can then easily cross-check that \begin{equation}
2\pi\int_{-1}^{+1}d\xi_{0}\int_{0}^{+\infty}p^{2}f_{M}\left(\psi,p,\xi\right)dp=\frac{2\pi p_{th}^{\dagger3}}{\beta_{th}^{\dagger3}}\int_{-1}^{+1}d\xi_{0}\int_{0}^{+\infty}\overline{p}^{2}\overline{f}_{M}\left(\psi,\overline{p},\xi\right)d\overline{p}=1\end{equation}
is well recovered$.$


\section{The non-maxwellian tail}

The tail part of the distribution function can be expressed in the
simple form

\begin{eqnarray}
f_{3T}\left(\psi,p,\xi\right) & = & \alpha_{3T}H\left(p_{max}-p\right)H\left(p-p_{min}\right)\times\nonumber \\
 &  & \left\{ \exp\left[-\frac{p_{\parallel}^{2}}{2\Theta_{\parallel f}\left(\psi\right)}\right]\exp\left[-\frac{p_{\perp}^{2}}{2\Theta_{\perp}\left(\psi\right)}\right]H\left(p_{\parallel}\right)\right.\nonumber \\
 &  & +\left.\exp\left[-\frac{p_{\parallel}^{2}}{2\Theta_{\parallel b}\left(\psi\right)}\right]\exp\left[-\frac{p_{\perp}^{2}}{2\Theta_{\perp}\left(\psi\right)}\right]\left(1-H\left(p_{\parallel}\right)\right)\right\} \label{eq:}\end{eqnarray}
\[
\]
where $H\left(x\right)$ is the usual Heaviside function while $\Theta_{\parallel f}$,
$\Theta_{\parallel b}$ and $\Theta_{\bot}$ are respectively the
parallel forward, backward and perpendicular pseudo-temperatures normalized
to $m_{e}c^{2}$. They are considered as pseudo-temperatures, since
they do not correspond to a thermodynamic equilibrium for which the
electron temperature $T_{e}$ has a true physical meaning, i.e. the
mean thermal velocity. They are introduced as simple parameters to
characterize the lack of symmetry around a given direction. Here,
the parallel component is expressed as\begin{equation}
p_{\Vert}=p\xi=\mathbf{p}\cdot\hat{b}\label{eq:}\end{equation}
where $\hat{b}$ is the unitary vector on the magnetic field line.
Conversely, $p_{\bot}$ is defined as $p^{2}=p_{\Vert}^{2}+p_{\bot}^{2}$.
The distribution function $f_{3T}$ is also normalized to unity. Let
$p_{min}$ being the momentum value corresponding to the intersection
of $f_{3T}$ with $f_{M}$, and $p_{max},$ the upper limit above
which it is zero.\begin{equation}
\int_{-\infty}^{+\infty}dp_{\Vert}\int_{0}^{+\infty}f_{3T}\left(\psi,p,\xi\right)p_{\perp}dp_{\perp}=1\label{eq:}\end{equation}


The presence of the Heaviside functions makes the calculation of the
coefficient $\alpha_{3T}$ non trivial. In the limits $lim_{p_{min\rightarrow0}}f_{3T}$
and $lim_{p_{max\rightarrow+\infty}}f_{3T}$, the parallel and perpendicular
dynamics are decoupled, and an analytical expression may be derived.
Indeed, In this case,

\begin{eqnarray}
\alpha_{3T}\left(\int_{0}^{+\infty}\exp\left[-\frac{p_{\parallel}^{2}}{2\Theta_{\parallel f}}\right]dp_{\Vert}+\int_{-\infty}^{0}\exp\left[-\frac{p_{\parallel}^{2}}{2\Theta_{\parallel b}}\right]dp_{\Vert}\right)\nonumber \\
\times\int_{0}^{+\infty}\exp\left[-\frac{p_{\perp}^{2}}{2\Theta_{\perp}}\right]2\pi p_{\perp}dp_{\perp} & = & 1\label{eq:}\end{eqnarray}
and a simple integration gives\begin{equation}
\int_{0}^{+\infty}\exp\left[-\frac{p_{\perp}^{2}}{2\Theta_{\perp}}\right]2\pi p_{\perp}dp_{\perp}=2\pi\Theta_{\perp}\label{eq:}\end{equation}
while\begin{equation}
\int_{0}^{+\infty}\exp\left[-\frac{p_{\parallel}^{2}}{2\Theta_{\parallel f}}\right]dp_{\Vert}=\sqrt{\frac{\pi}{2}}\sqrt{\Theta_{\parallel f}}\label{eq:}\end{equation}


Therefore, it is straightforward to find that \begin{equation}
\alpha_{3T}=\frac{1}{\sqrt{2\pi^{3}}}\frac{1}{\Theta_{\perp}\left(\sqrt{\Theta_{\parallel f}}+\sqrt{\Theta_{\parallel b}}\right)}\label{eq:}\end{equation}
and for the Maxwellian case $\Theta_{\Vert f}=\Theta_{\Vert b}=\Theta_{\bot}=\Theta$,
one recovers well the relation\begin{equation}
\alpha_{3T}=\alpha_{M}=\left(2\pi\Theta\right)^{-3/2}\label{eq:}\end{equation}


Considering the reference to the thermal value $p_{th}^{\dagger}$
, defining $\overline{p}_{\parallel}^{2}=\overline{p}_{\parallel}^{}p_{th}^{\dagger}$
and $\overline{p}_{\bot}^{2}=\overline{p}_{\bot}^{}p_{th}^{\dagger}$,
as well as for the temperatures, one finds \begin{eqnarray}
f_{3T}\left(\psi,\overline{p},\xi\right) & = & \overline{\alpha}_{3T}H\left(\overline{p}_{max}-\overline{p}\right)H\left(\overline{p}-\overline{p}_{min}\right)\times\label{eq:}\\
 &  & \left\{ \exp\left[-\frac{\overline{p}_{\parallel}^{2}}{2\overline{T}_{\parallel f}\left(\psi\right)}\right]\exp\left[-\frac{\overline{p}_{\perp}^{2}}{2\overline{T}_{\perp}\left(\psi\right)}\right]H\left(\overline{p}_{\parallel}\right)\right.\nonumber \\
 &  & \left.+\exp\left[-\frac{\overline{p}_{\parallel}^{2}}{2\overline{T}_{\parallel b}\left(\psi\right)}\right]\exp\left[-\frac{\overline{p}_{\perp}^{2}}{2\overline{T}_{\perp}\left(\psi\right)}\right]\left(1-H\left(\overline{p}_{\parallel}\right)\right)\right\} \nonumber \end{eqnarray}
using the identity for the Heaviside function\begin{equation}
H\left(ax+b\right)=H\left(x+\frac{b}{a}\right)H\left(a\right)+H\left(-x-\frac{b}{a}\right)H\left(-a\right)\label{eq:}\end{equation}
and the fact that $p_{th}^{\dagger}>0$. The modified normalized coefficient
is immediatly\begin{equation}
\overline{\alpha}_{3T}=\frac{1}{\sqrt{2\pi^{3}}}\frac{1}{\overline{T}_{\perp}\left(\sqrt{\overline{T}_{\parallel f}}+\sqrt{\overline{T}_{\parallel b}}\right)}\label{eq:}\end{equation}



\section{Calculation of $\overline{p}_{min}$}

For a given value of the coefficient $\lambda$, it is necessary to
find the intersection $\overline{p}_{min}$ between the Maxwellian
bulk $\overline{f}_{M}$ and the tail described by $\overline{f}_{3T}$.
Considering the normalisation of $\overline{f}$ to unity as the constraint,
it turns out that $\lambda$ and $\overline{p}_{min}$ are linked
by the relation\begin{equation}
\left(1-\lambda\right)\overline{\alpha}_{M}\left(1-\delta_{M}\left(\overline{p}_{min}\right)\right)+\lambda\overline{\alpha}_{3T}\left(1-\delta_{3T}\left(\overline{p}_{min}\right)\right)=1\label{eq:normeq}\end{equation}
where $\delta_{M}$ and $ $$\delta_{3T}$ are by definition small
corrections, i.e. $\delta_{M}\ll1$ and $\delta_{3T}\ll1$, provided
$\lambda\ll1$. Indeed, in this case, $\overline{p}_{min}\gg p_{th}^{\dagger}$,
while the contribution of the region $\overline{p}\leq\overline{p}_{min}$
to the deviation of the normalisation coefficient is almost negligible. 

The solution of the Eq.\ref{eq:normeq} can be obtained iteratively,
using the iterative Newton or Newton-Raphson methods or simply using
the dichotomy techniques.

\[
\]

\end{document}
